Die Abbildung ist auf

\mavergleichskettedisp
{\vergleichskette
{ \R^3
}
{ =} { \R^2 \times \R
}
{ } {
}
{ } {
}
{ } {
}
}
{}{}{}
das Produkt einer
\definitionsverweis {ebenen Drehung}{}{}
mit der Identität. Daher liegt eine euklidische
\definitionsverweis {Isometrie}{}{}
des $\R^3$ vor. Es sei

\mavergleichskette
{\vergleichskette
{ (x,y,z)
}
{ \in }{  T
}
{  }{
}
{    }{
}
{   }{
}
}
{}{}{}
und

\mavergleichskettedisp
{\vergleichskette
{ \Psi_\alpha(x,y,z)
}
{ =} { ( u,v,z)
}
{ } {
}
{ } {
}
{ } {
}
}
{}{}{.}
Dann ist

\mavergleichskettealigndrucklinks
{\vergleichskettealigndrucklinks
{ \left(  \sqrt{u^2+v^2} -R  \right)^2 +z^2
}
{ =} { \left(  \sqrt{ \left(  x \cos \alpha - y \sin  \alpha   \right)^2+ \left(  x \sin \alpha + y \cos  \alpha  \right)^2} -R  \right)^2 +z^2
}
{ =} { \left(  \sqrt{  x^2 \cos^{ 2  } \alpha + y^2 \sin^{ 2  }  \alpha  -2xy \cos \alpha \sin  \alpha    + x^2 \sin^{ 2  } \alpha + y^2 \cos^{ 2  }  \alpha  +2xy \cos \alpha \sin  \alpha     } -R  \right)^2 +z^2
}
{ =} { \left(  \sqrt{  x^2   + y^2 } -R  \right)^2 +z^2
}
{ =} { r^2
}
}
{}{}{,}
also

\mavergleichskette
{\vergleichskette
{ \Psi_\alpha(x,y,z)
}
{ \in }{  T
}
{  }{
}
{    }{
}
{   }{
}
}
{}{}{.}
Es ist also
\maabbdisp {\Psi_\alpha} { T} {T
} {}
ein Diffeomorphismus und eine Isometrie, da $T$ die induzierte riemannsche Struktur trägt.

 [[Kategorie:Latexseite]]
